% Part: first-order-logic
% Chapter: proof-systems
% Section: introduction

\documentclass[../../../include/open-logic-section]{subfiles}

\begin{document}

\iftag{FOL}
      {\olfileid{fol}{prf}{int}}
      {\olfileid{pl}{prf}{int}}

\olsection{Inleiding}

Logica's hebben doorgaans zowel een semantiek als
!!a{derivation}ssysteem. De semantiek betreft begrippen als waarheid,
vervulbaarheid, geldigheid en semantisch gevolg. Het doel van
!!{derivation}ssystemen is een zuiver syntactische methode te bieden
om semantisch gevolg en geldigheid vast te stellen. Ze zijn zuiver
syntactisch in die zin dat !!a{derivation} in zo'n systeem een eindig
syntactisch object is, gewoonlijk een rij (of een andere eindige
rangschikking) van !!{sentence}s of !!{formula}s. Goede
!!{derivation}ssystemen hebben de eigenschap dat van elke gegeven rij
of rangschikking van !!{sentence}s of !!{formula}s mechanisch kan
worden nagegaan of zij ``correct'' is.

De eenvoudigste (en historisch eerste) !!{derivation}ssystemen voor
de eerste-ordelogica waren \emph{axiomatisch}. Een rij !!{formula}s
geldt in zo'n systeem als !!a{derivation} als elke afzonderlijke
!!{formula} daarin ofwel tot een vaste verzameling ``axioma's''
behoort, ofwel uit eerdere !!{formula}s in de rij volgt door een van
een vast aantal ``afleidingsregels''. Bovendien moet mechanisch kunnen
worden nagegaan of !!a{formula} een axioma is en of zij volgens een
afleidingsregel correct uit andere !!{formula}s volgt. Axiomatische
!!{derivation}ssystemen zijn eenvoudig te beschrijven en ook
metatheoretisch eenvoudig te hanteren. !!{derivation}s erin zijn echter
moeilijk te lezen en te begrijpen, en eveneens moeilijk op te stellen.

Andere !!{derivation}ssystemen zijn ontwikkeld om !!{derivation}s
gemakkelijker te kunnen construeren of om voltooide !!{derivation}s
gemakkelijker te kunnen begrijpen. Voorbeelden zijn natuurlijke
deductie, waarheidsbomen, ook wel tableaubewijzen genoemd, en de
sequentencalculus. Sommige !!{derivation}ssystemen zijn speciaal met
het oog op mechanisering ontworpen. De resolutiemethode is bijvoorbeeld
eenvoudig in software te implementeren, maar haar !!{derivation}s zijn
vrijwel niet te begrijpen. De meeste van deze andere
!!{derivation}ssystemen stellen !!{derivation}s voor als bomen van
!!{formula}s in plaats van als rijen. Daardoor is gemakkelijker te
zien welke delen van !!a{derivation} van welke andere delen afhangen.

Voor een gegeven logica, zoals de eerste-ordelogica, leggen de
verschillende !!{derivation}ssystemen dus op verschillende manieren vast
wat het voor !!a{sentence} betekent om een \emph{stelling} te zijn
en om uit andere zinnen !!{derivable} te zijn. Hoe dit ook wordt gedaan
(met axiomatische !!{derivation}s, natuurlijke deducties,
sequent!!{derivation}s, waarheidsbomen of resolutieweerleggingen), we
willen dat deze relaties overeenkomen met de semantische begrippen
geldigheid en semantisch gevolg. We schrijven $\Proves !A$ voor
``$!A$~is een stelling'' en ``$\Gamma \Proves !A$'' voor
``$!A$~is uit~$\Gamma$ !!{derivable}.'' Ongeacht de definitie van
$\Proves$ willen we dat dit teken met $\Entails$ overeenkomt, dus:
\begin{enumerate}
\item $\Proves !A$ dan en slechts dan als $\Entails !A$
\item $\Gamma \Proves !A$ dan en slechts dan als $\Gamma \Entails !A$
\end{enumerate}
De richting ``alleen als'' hierboven heet \emph{correctheid}.
!!^a{derivation}ssysteem is correct als !!{derivability} semantisch
gevolg (of geldigheid) garandeert. Elk deugdelijk
!!{derivation}ssysteem moet correct zijn; incorrecte
!!{derivation}ssystemen zijn volstrekt onbruikbaar. Het gehele doel van
!!a{derivation} is immers een syntactische garantie van geldigheid of
semantisch gevolg te bieden. We zullen de correctheid bewijzen van de
!!{derivation}ssystemen die we behandelen.

De omgekeerde richting, ``als'', is eveneens belangrijk en heet
\emph{volledigheid}. Een volledig !!{derivation}ssysteem is krachtig
genoeg om aan te tonen dat $!A$~een stelling is wanneer $!A$~geldig is,
en dat $\Gamma \Proves !A$ wanneer $\Gamma \Entails !A$.
Volledigheid is moeilijker vast te stellen en sommige logica's hebben
geen volledige !!{derivation}ssystemen. De eerste-ordelogica heeft die
wel. Kurt G\"odel bewees in zijn proefschrift uit 1929 als eerste de
volledigheid voor !!a{derivation}ssysteem voor de eerste-ordelogica.

Een ander begrip dat met !!{derivation}ssystemen samenhangt, is
\emph{consistentie}. Een verzameling !!{sentence}s heet inconsistent
als om het even wat eruit kan worden !!{derive}d, en anders consistent.
Inconsistentie is de syntactische tegenhanger van onvervulbaarheid:
evenals onvervulbare verzamelingen vormen inconsistente verzamelingen
!!{sentence}s geen goede theorieën; ze vertonen een fundamenteel
gebrek. Consistente verzamelingen !!{sentence}s hoeven niet waar of
bruikbaar te zijn, maar voldoen ten minste aan deze minimale eis van
logische bruikbaarheid. De precieze definitie van consistentie van
verzamelingen !!{sentence}s kan per !!{derivation}ssysteem verschillen.
Net als bij~$\Proves$ willen we echter dat consistentie samenvalt met
haar semantische tegenhanger, vervulbaarheid. We willen dat $\Gamma$
altijd consistent is dan en slechts dan als zij vervulbaar is. De
richting ``alleen als'' komt hier neer op volledigheid (consistentie
garandeert vervulbaarheid) en de richting ``als'' op correctheid
(vervulbaarheid garandeert consistentie). Voor de klassieke
eerste-ordelogica zijn de twee versies van correctheid en volledigheid
inderdaad equivalent.
\end{document}
